\begin{Exercise}[label={ex:perfetti}]
    Scrivere un programma che chiede all’utente di inserire un numero naturale. Il programma verifica che il numero sia perfetto.
    Un numero si dice perfetto quando è uguale all somma di tutti i suoi divisori (eccetto se stesso).

    Scrivere le funzioni:
    \begin{itemize}
        \item leggiNumero(): che legge e restituisce un numero naturale
        \item numeroPerfetto(): che riceve in ingresso un numero naturale e restituisce vero
        \item sommaDivisori(): che riceve in ingresso un numero naturale e restituisce la somma dei suoi divisori (eccetto se stesso)
    \end{itemize}
\end{Exercise}

\begin{Answer}[ref={ex:perfetti}]
  	\lstinputlisting[style=cstylesmall]{src/perfetti.c}
\end{Answer}